\documentclass{situationist-preprint}

\title{Economics as an Idealized Science}
\msid{SCI-14-a}
\msfamily{SCI-14}
\msdiscipline{Philosophy of Science}
\mstype{Principal manuscript}
\mskeywords{economics; idealization; rational choice; models; approximation; experimental economics; demarcation}

\begin{document}
\maketitle

\begin{abstract}
Is economics a science? Its critics point to the falsity of its assumptions about agents and the weakness of its predictions. We argue that economics is a science, and an idealized one: its models are false of actual agents and apply to them only approximately. Economics satisfies the criterion of SCI-01: it has the institutions of review, replication and public criticism, and its record has been corrected by them. Its models are idealized in the sense of SCI-10. Whether a model applies approximately to a given phenomenon depends on whether its idealizations concern difference-makers for that phenomenon, which is an empirical question. Experimental markets show the competitive model applying to price convergence with traders far from rational, while expected-utility theory fails for individual choice under risk.
\end{abstract}

\section{The question}

\citet{Mill1844} held that political economy reasons from an abstraction, an agent who desires wealth and is constrained only by the means of obtaining it, and that its conclusions are true only in the abstract, as statements of tendencies that other causes may counteract. \citet{Friedman1953} argued that the realism of assumptions is irrelevant and that a theory is to be judged only by its predictions. \citet{Hausman1992} reconstructed economics as an inexact and separate science whose laws hold with ceteris paribus qualifications, and \citet{Rosenberg1992} argued that its failure to improve its predictions shows that it is not an empirical science. Recent work treats models as isolations of causal factors \citep{Maki2009} or as credible worlds \citep{Sugden2000}.

We argue that economics is a science, by the criterion of institutionalized error-correction defended in SCI-01, and an idealized one, in the sense of SCI-10: its models are false, and represent through partial similarity in difference-making respects.

\section{Economics is a science}

By the criterion of SCI-01, the question is whether economics subjects its accepted results to standing checks: review at admission, replication, and public criticism without terminal date. It does. Its journals are refereed, leading journals require deposit of data and code, and its results have been replicated on a large scale. When eighteen laboratory experiments published in leading economics journals were replicated, eleven reproduced the original effect \citep{Camerer2016}, a rate higher than in psychology but far from complete. The record has been corrected by these institutions. \citet{Leamer1983} showed that econometric results often depended on undisclosed choices of specification, and the subsequent turn to experimental and quasi-experimental designs corrected the field's practice \citep{AngristPischke2010}. The correction was slow, as SCI-19 leads one to expect, but it occurred. The usual objections, about false assumptions and weak predictions, do not bear on this criterion.

\section{Economics is idealized}

The models of economics assign to agents properties that no actual agent has: complete and transitive preferences over all options, unlimited computational capacity, and, in many models, perfect information. They assign to markets properties that no actual market has: infinitely many traders, homogeneous goods and costless transactions. By SCI-10, the models are false of every actual agent and market.

\begin{definition}[Approximate application]\label{d:approx}
A model \emph{applies approximately} to a class of targets for a phenomenon if it agrees with the targets, within a tolerance, in the respects that make a difference to the phenomenon.
\end{definition}

\begin{proposition}\label{p:range}
A model of economics applies approximately to actual agents for a phenomenon if and only if its idealizations concern features that make no difference to that phenomenon, or differences that remain within the tolerance.
\end{proposition}

\begin{proof}
By Definition~\ref{d:approx}, approximate application requires agreement only in the difference-making respects, and the idealizations are the respects of disagreement.
\end{proof}

Proposition~\ref{p:range} makes the range of application an empirical question, which experiments have answered in particular cases.

\paragraph{Market convergence.} \citet{Smith1962} showed that in double-auction markets with a handful of traders, each with private information about their own values and costs, prices converge quickly to the competitive equilibrium. The competitive model assumes infinitely many fully informed traders. \citet{GodeSunder1993} showed that convergence and high allocative efficiency are achieved even by computer traders who bid at random subject only to a budget constraint. The traders' rationality is not a difference-maker for convergence in that institution; the market rules are. The competitive model applies approximately, though false of every trader.

\paragraph{Choice under risk.} Expected-utility theory models agents whose preferences over lotteries are linear in probabilities. \citet{KahnemanTversky1979} showed systematic violations: people overweight small probabilities and evaluate outcomes as gains and losses relative to a reference point. For individual choice under risk the shape of the probability weighting is a difference-maker, and the model does not apply approximately; for aggregate phenomena in markets that discipline or average individual departures, it may.

\section{Objections and replies}

\begin{objection}[Friedman]
The realism of assumptions is irrelevant; only predictions matter \citep{Friedman1953}.
\end{objection}

\begin{reply}
Friedman was right that false assumptions do not disqualify a model, and wrong that their falsity is irrelevant. By Proposition~\ref{p:range}, whether a model applies depends on whether its false assumptions concern difference-makers, as they sometimes do.
\end{reply}

\begin{objection}[Rosenberg]
Economics has not improved its predictive power in two centuries, and so is not an empirical science.
\end{objection}

\begin{reply}
Quantitative prediction of open systems is poor in meteorology beyond some days and in geology always, and neither is thereby unscientific. Economics has improved its knowledge of which phenomena its models apply to, which is the improvement Proposition~\ref{p:range} predicts for an idealized science.
\end{reply}

\begin{objection}[Behavioral economics]
Behavioral economics shows that the core assumptions are false: economics is refuted, not idealized.
\end{objection}

\begin{reply}
The findings show that the assumptions are false, which an idealized science grants from the start, and they mark out the phenomena for which the assumptions are difference-makers. Their incorporation into the field is itself the error-correction that makes economics a science.
\end{reply}

\section{Conclusion}

Economics has the institutions of error-correction that distinguish science. Its models are false of every actual agent and apply approximately where their idealizations make no difference to the phenomenon. Economics is a science, and an idealized one.

\begin{thebibliography}{99}\small
\bibitem[Angrist and Pischke(2010)]{AngristPischke2010} Angrist, J.~D. and Pischke, J.-S. (2010). The credibility revolution in empirical economics: How better research design is taking the con out of econometrics. \emph{Journal of Economic Perspectives} 24(2): 3--30.
\bibitem[Camerer et~al.(2016)]{Camerer2016} Camerer, C.~F., et~al. (2016). Evaluating replicability of laboratory experiments in economics. \emph{Science} 351: 1433--1436.
\bibitem[Friedman(1953)]{Friedman1953} Friedman, M. (1953). The methodology of positive economics. In \emph{Essays in Positive Economics}. University of Chicago Press.
\bibitem[Gode and Sunder(1993)]{GodeSunder1993} Gode, D.~K. and Sunder, S. (1993). Allocative efficiency of markets with zero-intelligence traders: Market as a partial substitute for individual rationality. \emph{Journal of Political Economy} 101: 119--137.
\bibitem[Hausman(1992)]{Hausman1992} Hausman, D.~M. (1992). \emph{The Inexact and Separate Science of Economics}. Cambridge University Press.
\bibitem[Kahneman and Tversky(1979)]{KahnemanTversky1979} Kahneman, D. and Tversky, A. (1979). Prospect theory: An analysis of decision under risk. \emph{Econometrica} 47: 263--291.
\bibitem[Leamer(1983)]{Leamer1983} Leamer, E.~E. (1983). Let's take the con out of econometrics. \emph{American Economic Review} 73: 31--43.
\bibitem[M\"aki(2009)]{Maki2009} M\"aki, U. (2009). MISSing the world: Models as isolations and credible surrogate systems. \emph{Erkenntnis} 70: 29--43.
\bibitem[Mill(1844)]{Mill1844} Mill, J.~S. (1844). On the definition of political economy; and on the method of investigation proper to it. In \emph{Essays on Some Unsettled Questions of Political Economy}. Parker. (First published 1836.)
\bibitem[Rosenberg(1992)]{Rosenberg1992} Rosenberg, A. (1992). \emph{Economics---Mathematical Politics or Science of Diminishing Returns?} University of Chicago Press.
\bibitem[Smith(1962)]{Smith1962} Smith, V.~L. (1962). An experimental study of competitive market behavior. \emph{Journal of Political Economy} 70: 111--137.
\bibitem[Sugden(2000)]{Sugden2000} Sugden, R. (2000). Credible worlds: The status of theoretical models in economics. \emph{Journal of Economic Methodology} 7: 1--31.
\end{thebibliography}

\end{document}
